\documentclass[10pt,a4paper]{amsart}
\usepackage[margin=19mm]{geometry}
\usepackage{amsmath,amssymb,mathtools}
\usepackage[scr=rsfs,cal=euler]{mathalfa}
\usepackage{xfrac}
\usepackage[unicode,hidelinks,bookmarks=false]{hyperref}
\newcommand{\ie}{i.e.\ }
\newcommand{\cf}{cf.\ }
\newcommand{\resp}{respectively\ }
\newcommand{\Subsection}{\subsection}
\providecommand{\nobreakdash}{\nobreak\hbox{-}}
\setlength{\emergencystretch}{2em}
\multlinegap0pt
\usepackage{bm}
\usepackage{bbm}
\usepackage{dsfont}
\usepackage{xspace}
\usepackage{cancel}
\usepackage{mathtools}
\usepackage{amsthm}
\makeatletter
\@ifundefined{lemma}{\newtheorem{lemma}{Lemma}}{}
\makeatother
\DeclareRobustCommand{\stirling}{\genfrac\{\}{0pt}{}}
\title[On a problem on a generalization of Euler's totient function]{On a problem on a generalization of Euler's totient function}
\author{John M. Campbell}
\date{Source: arXiv:2606.01633v1 (CC BY 4.0)}
\begin{document}
\maketitle

\noindent\emph{Source and licence.} On a problem on a generalization of Euler's totient function. Version arXiv:2606.01633v1 (CC BY 4.0), distributed under the Creative Commons Attribution 4.0 International licence, as recorded in the arXiv licence metadata for this exact version and shown on the abstract page https://arxiv.org/abs/2606.01633v1. Full rights notice: licenses/arxiv-2606.01633-CC-BY-4.0.txt.\par\smallskip

\noindent\textit{Editorial scope.} Counted unit: the Lemma whose source label is \texttt{triangularlemma} in the article (source0.tex lines 122--126, source label \texttt{triangularlemma}), delivered together with the article's own complete original proof of it (source0.tex lines 128--233, one \texttt{proof} environment of 106 lines). No mathematics is written, repaired or re-derived: statement and proof are the article's own text line for line. Required context retained uncounted: defineBNx (lines 110--112); DStirling (lines 116--118). Every definition, display and label of those lines is retained unchanged. Omitted from this package and not redistributed: 1--109, 113--115, 119--121, 127--127, 234--526. Nothing inside a retained excerpt is omitted or altered: every retained body is delivered exactly as the article's lines are written, so no omission receipt of any kind accompanies this package. Citations: the retained excerpts carry no citation command at all, so no bibliography entry of the article is transcribed and no reference is redistributed. Reference adaptation: none was needed and none was made; every reference inside the delivered unit resolves inside it, so no unresolved reference or citation is produced. Printed slips kept literal: no printed word, symbol, hypothesis or number of the counted statement or of its proof was altered. Layout and class: the article's own class is \texttt{article} with packages including amssymb, amsmath, amsthm, hyperref, graphicx, xcolor, mathrsfs; the wrapper is the lane's pinned \texttt{amsart} editorial wrapper at 10pt on A4, which restates the theorem-like environments the retained text needs and adds the article's own macro definitions that the retained text needs (\texttt{\textbackslash stirling}), with extra wrapper packages (bm, bbm, dsfont, xspace, cancel, mathtools). The delivered numbering is the unit's own. Assets: no retained span contains a figure environment, a graphics-inclusion command, a PDF-page inclusion, a table or a TikZ picture; no image file is copied into this package, the package declares no asset dependency and no image file is redistributed with it. Licence: version arXiv:2606.01633v1 of this article is distributed under the Creative Commons Attribution 4.0 International licence, as recorded in the arXiv licence metadata for this exact version and shown on the abstract page https://arxiv.org/abs/2606.01633v1. A full rights notice is delivered at licenses/arxiv-2606.01633-CC-BY-4.0.txt. Characters: every retained excerpt is pure ASCII, so the delivered engine, XeLaTeX with its default Latin Modern text font, reports no missing character.
\begin{equation}\label{defineBNx} 
 B_{N}(x) = \sum_{r=0}^{N} \binom{N}{r} B_{r} x^{N-r}
\end{equation}

\begin{equation}\label{DStirling} 
 D^{m} = \sum_{\ell = 1}^{m} \stirling{m}{\ell} x^{\ell} \frac{d^{\ell}}{dx^{\ell}} 
\end{equation}

\begin{lemma}\label{triangularlemma}
 For $k > 1$ odd, suppose that $C_k(v)$ is an integer for every $v \geq 1$. Then, for every even integer $h \in [2, k-1]$, 
 we have that 
 $ 2 \frac{k!}{h!} B_{h} \in \mathbb{Z}$. 
\end{lemma}

\begin{proof}
 Write $N = k + 1$, with $N \geq 4$ being even. For an integer 
\begin{equation}\label{srange} 
 s \in \left[ 1, \frac{N-2}{2} \right], 
\end{equation}
 we define
\begin{equation}\label{defineTs}
 T_{N, s} = T_{s} = 2 \frac{(N-1)!}{(N-2s)!} B_{N-2s}. 
\end{equation}
 Since every even integer $h \in [2, k-1]$ can be written uniquely in the form
 $h = N - 2s$ for $N$ and $s$ as specified, 
 it remains to prove that $T_{s} \in \mathbb{Z}$. 

 Now, define 
\begin{equation}\label{defineFN}
 F_{N}(x) = x^{1-N} \big( B_{N}(x) - B_{N} \big), 
\end{equation}
 recalling the definition of Bernoulli polynomials in \eqref{defineBNx}. We claim that 
\begin{equation}\label{nestedD}
 C_{N - 1}(2s) = \frac{2}{N} \big( D^{2s} F_{N} \big)(1). 
\end{equation}
 To prove the relation in \eqref{nestedD}, we begin by manipulating the formula in 
 \eqref{defineBNx} to obtain that 
\begin{equation}\label{manipulatedBNx} 
 x^{1-N} B_{N}(x) = \sum_{r=0}^{N} \binom{N}{r} B_{r} x^{1-r},
\end{equation}
 so that \eqref{manipulatedBNx} gives us that 
\begin{equation}\label{sumtoreduce} 
 D^{2s}\big( x^{1-N} B_{N}(x) \big) \Big|_{x=1}
 = \sum_{r=0}^{N} \binom{N}{r} B_{r}(1-r)^{2s}. 
\end{equation}
 From the vanishing of the $r = 1$ case of the summand on the right of \eqref{sumtoreduce} together with the vanishing of Bernoulli 
 numbers with odd indices greater than $1$, we find that 
\begin{equation}\label{combinewithFN} 
 D^{2s}\big( x^{1-N} B_{N}(x) \big) \Big|_{x=1}
 = \frac{N}{2} C_{N-1}(2s) + B_{N}(1-N)^{2s}, 
\end{equation}
 so that an application of the definition of $F_{N}(x)$ in \eqref{defineFN}
 via \eqref{combinewithFN} gives us an equivalent version of \eqref{nestedD}. 

 Now, we rewrite $F_{N}(x)$, as defined in \eqref{defineFN}, by setting $g(x) = x^{1-N}$ and $H(x) = B_{N}(x) - B_{N}$, with $F_{N}(x) 
 = g(x) H(x)$. Observe that $B_{N}(1) = B_{N}$ and $H(1) = 0$.
 The product rule for derivatives then gives us that 
\begin{equation}\label{diffprod} 
 F_{N}^{(\ell)}(1) = \sum_{r=1}^{\ell} \binom{\ell}{r} g^{(\ell - r)}(1) B_{N}^{(r)}(1),
\end{equation}
 noting that we obtain the vanishing of $\binom{\ell}{r} g^{(\ell - r)}(1) H^{(r)}(1) $
 for the $r = 0$ case. 
 Observe that $g^{(\ell - r)}(1) \in \mathbb{Z}$, recalling the definition of $g(x)$ as an integer power of $x$. 
 From the definition of Bernoulli polynomials in \eqref{defineBNx}, one may verify that 
\begin{equation}\label{setx1} 
 B_{N}^{(r)}(x) = \frac{N!}{(N-r)!} B_{N - r}(x) 
\end{equation}
 for positive integers $r \leq N$, and we specialize \eqref{setx1} to 
\begin{equation}\label{specializeto1}
 B_{N}^{(r)}(1) = \frac{N!}{(N-r)!} B_{N - r}(1). 
\end{equation}
 According to the product rule expansion in \eqref{diffprod}, we have that sum in this product rule is restricted to indices $r \in [1, 
 \ell]$. In view of the identity in \eqref{nestedD}, which involves the application of $D^{2s}$ to $F_{N}(1)$, it is implicit that $\ell \leq 2 
 s$, and we recall the bounds on $s$ in \eqref{srange}, giving us that $2 \leq N - r$. Now, if $r$ is odd, then $N - r \geq 2$ is odd. 
 So, if $r$ is odd, then $B_{N-r}(1) = B_{N-r} = 0$. So, from \eqref{specializeto1}, we have that if $r$ is odd, 
 then $B_{N}^{(r)}(1)$ vanishes. 
 So, for an even value $r = 2i$, 
 we find that 
\begin{equation}\label{BntoT} 
 \frac{2}{N} B_{N}^{(2i)}(1) = T_{i}, 
\end{equation}
 recalling the definition in \eqref{defineTs}. An application of 
 the expression in \eqref{nestedD} of $C_{N-1}(2s)$ in terms of $\big( D^{2s}F_{N} \big)(1)$ 
 together with 
 the classical Stirling number identity in \eqref{DStirling}
 give  us that 
\begin{equation}\label{separateStirling} 
 \big( D^{2s} F_{N} \big)(1) = F_{N}^{(2s)}(1)
 + \sum_{\ell = 1}^{2s-1} \stirling{2s}{\ell} F_{N}^{(\ell)}(1). 
\end{equation}
 Through an application of \eqref{diffprod} to \eqref{separateStirling}, we obtain 
\begin{multline*}
 \frac{2}{N} \big( D^{2s} F_{N} \big)(1) = \sum_{r=1}^{2s} \binom{2s}{r} g^{(2s - r)}(1) \frac{2}{N} B_{N}^{(r)}(1) + \\ 
 \sum_{\ell = 1}^{2s-1} \stirling{2s}{\ell} \sum_{r=1}^{\ell} \binom{\ell}{r} g^{(\ell - r)}(1) \frac{2}{N} B_{N}^{(r)}(1), 
\end{multline*}
 and we rewrite this as 
\begin{multline*}
 \frac{2}{N} \big( D^{2s} F_{N} \big)(1) = \frac{2}{N} B_{N}^{(2s)}(1) + \sum_{r=1}^{2s-1} \binom{2s}{r} g^{(2s - r)}(1) \frac{2}{N} B_{N}^{(r)}(1) 
 + \\ 
 \sum_{\ell = 1}^{2s-1} \stirling{2s}{\ell} \sum_{r=1}^{\ell} \binom{\ell}{r} g^{(\ell - r)}(1) \frac{2}{N} B_{N}^{(r)}(1). 
\end{multline*}
 For indices $r$ involved in the above sums, we have that $1 \leq r \leq 2s \leq N-2$, 
 with $N - r \geq 2$. So, if $r$ is odd, then $N - r$ is odd and $\geq 3$, 
 with $B_{N-r}(1) = B_{N - r} = 0$, and hence
\begin{multline*}
 \frac{2}{N} \big( D^{2s} F_{N} \big)(1) = \frac{2}{N} B_{N}^{(2s)}(1) + \sum_{r=1}^{s-1} \binom{2s}{2r} g^{(2s - 2r)}(1) \frac{2}{N} B_{N}^{(2r)}(1) 
 + \\ 
 \sum_{\ell = 1}^{2s-1} \stirling{2s}{\ell} \sum_{ r = 1 }^{ \lfloor {\ell}/{2} \rfloor} \binom{\ell}{2r} g^{(\ell - 2r)}(1) \frac{2}{N} B_{N}^{(2r)}(1). 
\end{multline*}
 From the formula for $C_{N - 1}(2s)$ in \eqref{nestedD} together with the definition in \eqref{BntoT}, we find that there exist integers 
 $a_{s, r}$ such that 
\begin{equation}\label{inductionT} 
 T_{s} = C_{N-1}(2s) - \sum_{r=1}^{s-1} a_{s, r} T_{r}. 
\end{equation}
 For the $s = 1$ case of \eqref{inductionT}, we find that $C_{k}(2) = T_{1}$, and $C_{k}(2) $ is an integer by assumption. So, from the 
 assumption that $C_{k}(v)$ is an integer for every $v \geq 1$ and for $k > 1$ odd, and since $N \geq 4$ is even,  an inductive argument 
 gives us, from  \eqref{inductionT}, that $T_{s}$ is an integer for 
 $s$ in the interval indicated in \eqref{srange}, 
 as desired. 
\end{proof}

\end{document}
